\begin{pastexam}[label={px:ust-108-ee-ds-02}]{臺聯大電機類 108 學年度 第 2 題}
（問答）

(20\%) Given a link list structure as
\begin{minted}[fontsize=\small]{c}
struct Node
{
    int data;
    struct Node* next;
};
\end{minted}
A list head: \texttt{struct Node *Head}, A particular middle list element \texttt{struct Node *Mid}.
Write a program to delete \texttt{Mid} from the list.
\end{pastexam}

\begin{pastsoln}
單向串列的節點只知道後面是誰，不知道前面是誰。 要把 \texttt{Mid} 拿掉，得改的是\textbf{它前一個
節點}的 \texttt{next}，所以要從 \texttt{Head} 往後找，找到 \texttt{next} 等於 \texttt{Mid} 的節點
\texttt{prev}，再讓 \texttt{prev->next} 跳過 \texttt{Mid}（資結書虛擬碼~5.3）。

題目說 \texttt{Mid} 是中間的元素，照理不會是第一個；下面的程式還是把「\texttt{Mid} 就是第一個」
和「\texttt{Mid} 不在串列裡」也處理掉，因為第一個節點被刪時 \texttt{Head} 本身要改，所以參數用
\texttt{struct Node **head}。

\begin{minted}[fontsize=\small]{c}
#include <stdlib.h>

void delete_node(struct Node **head, struct Node *mid)
{
    struct Node *prev;

    if (*head == NULL || mid == NULL)
        return;
    if (*head == mid) {                 /* 刪第一個：Head 往後移 */
        *head = mid->next;
        free(mid);
        return;
    }
    prev = *head;                       /* 找 mid 的前一個節點 */
    while (prev->next != NULL && prev->next != mid)
        prev = prev->next;
    if (prev->next == mid) {            /* 找到了才刪 */
        prev->next = mid->next;
        free(mid);
    }
}
\end{minted}
呼叫時寫 \texttt{delete\_node(\&Head, Mid);}。

\textbf{正確性。} 迴圈不變量：\texttt{prev} 是串列中某個節點，而且從 \texttt{*head} 到
\texttt{prev} 的節點都不是 \texttt{mid}。 一開始 \texttt{prev = *head}，前面已經確定
\texttt{*head != mid}，成立；每次往後移一格之前都檢查過 \texttt{prev->next != mid}，所以仍成立。
迴圈結束時不是 \texttt{prev->next == mid}（找到前一個節點），就是走到最後一個節點都沒找到。 找到時，
\texttt{prev->next = mid->next} 讓串列從 \texttt{prev} 直接接到 \texttt{mid} 後面那個節點（\texttt{mid}
是最後一個時就是 \texttt{NULL}），其他節點的順序不變；最後才 \texttt{free(mid)}，先後不能顛倒，否則
\texttt{mid->next} 會讀到已經釋放的記憶體。

例如串列 $10 \to 20 \to 30 \to 40 \to 50$，\texttt{Mid} 指向 $30$：\texttt{prev} 從 $10$ 走到 $20$
（$20$ 的 \texttt{next} 就是 \texttt{Mid}），改成 $20 \to 40$，結果是 $10 \to 20 \to 40 \to 50$。

\textbf{複雜度。} 最壞要從頭走到 \texttt{Mid} 前面，$\Theta(n)$ 時間，$\Theta(1)$ 額外空間。
真正改指標只要 $\Theta(1)$，花時間的是找前一個節點。

\textbf{另一種寫法（只能用在 \texttt{Mid} 不是最後一個時）。} 不找前一個節點，而是把後一個節點
的內容搬到 \texttt{Mid}，再刪掉後一個：
\begin{minted}[fontsize=\small]{c}
void delete_middle(struct Node *mid)   /* mid 必須有下一個節點 */
{
    struct Node *q = mid->next;
    mid->data = q->data;
    mid->next = q->next;
    free(q);
}
\end{minted}
這是 $\Theta(1)$，題目說 \texttt{Mid} 是中間的元素，剛好可以用。 但它有兩個限制，作答時要寫出來：
\texttt{Mid} 是最後一個節點時沒有後一個可以搬，這招不能用；而且實際被釋放的是後一個節點，如果程式
別的地方還留著指向那個節點的指標，就會指到已經釋放的記憶體。
\end{pastsoln}
